\title{Kuznetsov components of Fano varieties}
\author{\textit{Based on a Lecture series by} Arend Bayer.}
\date{December 2023}
\maketitle

\newcommand{\deri}{\mathop{\mathbf{D}}}
\newcommand{\Jac}{\mathop{\mathcal{J\kern -0.5pt a\kern -0.5pt c\kern -1pt}}}
\newcommand{\IJac}{\mathop{I\mathcal{J\kern -0.5pt a\kern -0.5pt c\kern -1pt}}}
\newcommand{\Moduli}{\ensuremath{\mathscr{M}}}

We will discuss methods to study Kuznetsov components of derived categories of
coherent sheaves, using moduli theory and stability spaces. Along the way we
will look at classical questions regarding invariants and techniques surrounding
Fano varieties.

\subsection{About} This course on Fano varieties was a part of a Winter
School for early career researchers, organised by the UK Algebraic
Geometry Network. The lectures were delivered in-person in the University of
Warwick, and were transcribed by Parth Shimpi. These notes have undergone
several amendments and are not a verbatim recall of the lectures, therefore
discretion is advised when using this material. They are available online at
\url{https://pas201.user.srcf.net/documents/2023-ukag-kuznetsov.pdf}.

All errors and corrections should be communicated to by email to
\url{pas201@srcf.net}.


\section{Low--dimensional Fano varieties}
\begin{definition}
    A smooth variety \(X/\mathbb{C}\) is called \textit{Fano} if the
    anticanonical bundle \(\mathscr{O}(-K_X)=\Exter^\text{top}T_X\) is ample.
\end{definition}

Thus the only one-dimensional Fano variety is \(\mathbb{P}^1\).

More generally, all projective spaces and complete intersections of low degree
in projective spaces are Fano. In particular, any complete intersection
\(X\subset \mathbb{P}^{n+N}\) cut out by equations of degrees \(d_1,...,d_N\) is
Fano if and only if \(\sum_i d_i\leq n+N\). For example, this can be used to
show that all Fano 2-folds (i.e.\ del Pezzo surfaces) arise as \(\mathbb{P}^2\)
blown up at \(\leq 8\) points, or as the smooth quadric \(\mathbb{P}^1\times
\mathbb{P}^1\).

\begin{example}[Grassmannians]
    Generalising the example of projective space, each Grassmannian variety
    \(G=\text{Gr}(n,m)\) parametrising \(n\)-dimensional subspaces \(U_n\subset
    \mathbb{C}^m\) is Fano. To see why, we will compute its tangent bundle. Note
    we have the tautological exact sequence of vector bundles
    \(\mathcal{U}_n\into \mathscr{O}_X^{\oplus m}\onto \mathcal{Q}_{m-n}\), and
    that at a point \(p\in X\) corresponding to \(U_n\into \mathbb{C}^m\onto
    Q_{m-n}\) the tangent space \(T_pX\) is given by elements of \(\Hom(U_n,
    Q_{m-n})\) i.e.\ deformations of the subspace \(U_n\into \mathbb{C}^m\) up
    to automorphisms of \(U_n\). Thus the tangent bundle is
    \(T_X=\Hom(\mathcal{U}_n, \mathcal{Q}_{m-n})\), and we have
    \(\Exter^\text{top} T_X= (\Exter^\text{top} \mathcal{U}_n\dual)^{\otimes
    (m-n)}\otimes (\Exter^\text{top}\mathcal{Q}_{m-n})^{\otimes n}\). But the
    tautological exact sequence shows \(\Exter^\text{top}\mathcal{U}_n =
    (\Exter^\text{top}\mathcal{Q}_{m-n})^{-1}\), and the Pl\"ucker embedding
    \(G\into \mathbb{P}(\Exter^n \mathbb{C}^m)\) shows
    \(\Exter^\text{top}\mathcal{U}_n = \mathscr{O}(1)\) is ample. We conclude
    that so is \(\Exter^\text{top}T_X = \mathscr{O}_X(m)\).
\end{example}

Likewise, low degree subvarieties of Grassmannians are Fano, and this gives us
examples of Fano varieties in dimensions three and four. For example,
\(\text{Gr}(2,4)\subset \mathbb{P}^5\) is a smooth quadric 4-fold. Intersecting
\(\text{Gr}(2,5)\subset \mathbb{P}^9\) with three hyperplanes, we get the Fano 3-fold \(Y_5=\text{Gr}(2,5)\cap \mathbb{P}^6\). Likewise we have the
\textit{Gushel--Mukai 3-fold} \(X_{10}\) obtained by intersecting
\(\text{Gr}(2,5)\subset \mathbb{P}^9\) with two hyperplanes and a quadric.
Another similar example is \(X_{14}=\text{Gr}(2,6)\cap \mathbb{P}^9\subset
\mathbb{P}^{14}\).

\subsection{Deformation families}
Fano varieties of fixed dimension are \textit{bounded}, i.e.\ occur in
finitely many deformation families. In low dimensions, there are fairly explicit
unirational moduli spaces while the moduli spaces in higher dimensions are
studied using K-stability.

For Fano 3-folds, there are 105 families. Most have relatively low Picard rank--
this is because the behaviour is fairly complicated for high Picard ranks and in
fact sufficiently high Picard ranks simply become difficult to realise. For
instance the only way to obtain a Fano 3-fold of Picard rank \(8\) is the
trivial way, by taking the product of \(\mathbb{P}^1\) with a del Pezzo surface.

Thus we focus on the simplest case, that of Picard rank \(1\). Such
Fano 3-folds are called \textit{prime}, and have a canonical ample divisor \(H\)
(called the \textit{fundamental divisor}) which generates \(\Pic X\). We say the
\textit{index} of \(X\) is the integer \(i_X\) satisfying \(K_X+i_X \cdot H
\equiv 0\), and the \textit{degree} of \(X\) is the self-intersection
\(d_X=-K_X^3\). It is also convenient to introduce the number
\(g_X=\frac{1}{2}(d_X+2)\), called the \textit{genus} (the reason for this
terminology will be explained later.)

\begin{remark}
    If \(X\) is a prime Fano 3-fold, then we must have \(1\leq i_X\leq 4\).
    Indeed the lower bound comes from the Fano condition. To find the upper
    bound, note that if \(S\in |H|\) is a general element in the linear system
    then \(K_S=-(i_X-1)H|_S\) by adjunction. In particular if \(i_X>1\) then
    \(S\) is a del Pezzo surface and we have the well-known bound
    \[1\leq K_S^2=(i_X-1)^2H^3\leq 9.\]
    which shows \(i_X\) cannot exceed \(4\). Note \(i_X=4\) forces \(H^3=1\),
    and the only such Fano 3-fold is \(\mathbb{P}^3\).
\end{remark}

The deformation type of \(X\) is determined by \(i_X\) and \(d_X\), so it
suffices to compute the possible values of \(d_X\) for each \(i_X\in
\{1,2,3,4\}\). For \(i_X=3\) the only possibility is \(H^3=2\) i.e.\ \(d_X=54\)
(\(H^3=1\) is ruled out since the genus is an integer) and the variety is a
quadric in \(\mathbb{P}^4\).

For \(i_X=2\) the above remark shows \(1\leq d_X\leq 9\), but
in fact only \(1\leq d_X\leq 5\) are realised. These are given by \(Y_d\)
(\(1\leq d\leq 5\)), where \(Y_5\) was described above. \(Y_4\subset
\mathbb{P}^5\) is a \((2,2)\) complete intersection, \(Y_3\subset \mathbb{P}^4\)
is a cubic 3-fold, and \(Y_2\onto \mathbb{P}^3\) is a double cover.

For \(i_X=1\), we have the following result.

\begin{theorem}[(Ishkowskii)]
    If \(X\) is a prime Fano 3-fold of index 1, then we have \(2\leq g_X \leq
    12\) and \(g_X\neq 11\). This is called the \emph{Fano range}.
\end{theorem}

Thus there are \underline{seventeen} families of prime Fano 3-folds.

\subsection{A closer look at index 1} There is a beautiful correspondence
between Fano 3-folds of index \(1\), K3 surfaces, and curves which was first
shown by Mukai. Indeed in this case the anticanonical bundle \(-K_X=H\) is a
natural polarisation, and a general element in the linear system
\(|H|\) is a smooth K3 surface \(S\).

Moreover this surface naturally comes with the polarisation \(H|_S\), and a general element in this linear system is a genus \(g_S\) curve (the genus of a polarised K3 surface is defined this way). Then it can be shown that \(g_S\) coincides with \(g_X\).

Mukai showed that general K3 surfaces arise in this way.

\begin{theorem}[(Mukai)]
    If \(g\) is in the Fano range, then a general polarised K3 surface of genus
    \(g\) arises as a hyperplane section of a Fano 3-fold. Moreover, a general
    canonical curve arises as a hyperplane section of a polarised K3 (and hence
    as a complete intersection of a Fano 3-fold).
\end{theorem}

\subsection{Irrationality of the cubic 3-fold} A discussion on low dimensional
Fano varieties would not be complete without mentioning this famous result.
If \(X\) is a Fano 3-fold, then \(\HH^{3,0}(X)=\HH^0(X,\mathscr{O}_X(K_X))=0\)
and hence we can consider \(\HH^3(X,\mathbb{Z})\subset
\HH^3(X,\mathbb{C})=\HH^{2,1}(X)\oplus \HH^{1,2}(X)\). The space
\(\IJac X=\HH^{2,1}(X,\mathbb{C})\dual / \HH_3(X,\mathbb{Z})\) is a
polarised Abelian variety called the \textit{intermediate Jacobian of \(X\)}.

The behaviour of this under birational transformations is well-understood.
Firstly it is invariant under blowing up in codimension \(2\) locii, i.e.\ we
have \(\IJac(\text{Bl}_{\text{pts}}X)=\IJac X\). On the other hand if
\(C\subset X\) is a curve, then \(\IJac(\text{Bl}_CX)=\IJac X\times
\Jac C\). Since Jacobians of curves are irreducible as Abelian varieties and
\(\mathbb{P}^2\) has trivial intermediate Jacobian, we use weak factorisation to
deduce that \(\IJac X\) is a product of Jacobians whenever \(X\) is rational.
This is used to deduce the following.

\begin{theorem}[(Clemens--Griffiths)]
    A general cubic 3-fold is irrational.

    \begin{proof}[Sketch]
        If \(X\) is a cubic 3-fold, then \(\IJac X\) is a five dimensional
        principally polarised Abelian variety. We compute that the theta divisor
        \(\Theta\subset \text{IJac}X\) has minimal degree and an isolated
        singularity \(x_0\). This can happen only if \(\IJac X\) is
        indecomposable and not a Jacobian, and hence \(X\) is not rational.
    \end{proof}
\end{theorem}

Similar arguments also give Torelli theorem for cubic 3-folds, by
showing the projectivised tangent cone \(\mathbb{P}T_{x_0}\Theta\) is isomorphic
to \(X\). Thus \(X\) is determined by the principally polarised Abelian variety
\((\IJac X,\Theta)\), which is in turn determined by the polarised Hodge
structure on \(\HH^3(X,\mathbb{C})\).

\section{Semiorthogonal decompositions}
Given a variety \(X/\mathbb{C}\), recall the construction of the derived
category of coherent sheaves \(\deri(X)=\deri^b\Coh X\). If \(X\) is Fano, then
a classical result of Bondal--Orlov shows that \(\deri(X)\) determines \(X\). We
will refine this result and ask if \(X\) is determined by certain canonical
subcategories.

\begin{definition}
    An object \(E\in \deri(X)\) is called \textit{exceptional} if
    \(\Hom^{\bullet}(E,E)=\mathbb{C}[0]\). An \textit{exceptional collection} is
    an ordered sequence \((E_1,...,E_n)\) of exceptional objects such that
    \(\Hom^\bullet (E_i,E_j)=0\) whenever \(i>j\). We say an exceptional
    collection is \textit{full} if it generates \(\deri(X)\) (i.e.\ is not
    contained in a proper triangulated subcategory).
\end{definition}

Fano varieties are naturally equipped with exceptional collection. For example,
a classical result of Beilinson shows \(\mathbb{P}^n\) is has a full exceptional
collection \((\mathscr{O}_{\mathbb{P}^n}, \mathscr{O}_{\mathbb{P}^n}(1),...,
\mathscr{O}_{\mathbb{P}^n}(n))\). This generalises to the following result.

\begin{lemma}
    If \(X\) is Fano and \(L\in \deri(X)\) is a line bundle, then \(L\) is
    exceptional. Moreover, if \(L\) is an ample line bundle such that
    \(\mathscr{O}(-K_X)=L^{\otimes r}\) then
    \((\mathscr{O}_X, L,...,L^{\otimes r-1})\) is an
    exceptional collection.
    \begin{proof}
        We will show the second part of the statement, the first is analogous.
        We need to compute the group \(\Ext^q(L^{\otimes i},L^{\otimes j})\) for
        \(r> i\geq j \geq 0\). This is equal to \(\HH^q(X,
        \mathscr{O}_X(-K_X)\otimes L^{\otimes r+j-i})\). We know
        \(r+j-i>0\) and hence \(L^{\otimes r+j-i}\) is ample. Thus the required
        cohomology vanishes for \(q>0\) by the Kawamata--Veibig vanishing
        theorem. On the other hand if \(q=0\) and \(i>j\) then the group is
        \(\HH^0(X,L^{\otimes j-i})\) which vanishes since anti-ample bundles
        don't have global sections.
    \end{proof}
\end{lemma}

Weakening the definition of a full exceptional collection, we have the
following.

\begin{definition}
    If \(\mathscr{D}\) is a triangulated category, then a
    \textit{semi-orthogonal decomposition} denoted \(\mathscr{D}=\langle
    \mathscr{D}_1,\mathscr{D}_2,...,\mathscr{D}_n\rangle\) is a sequence of full
    triangulated subcategories that generate \(\mathscr{D}\) and satisfy
    \(\mathscr{D}_i\subset \mathscr{D}_j\orth\) whenever \(i<j\).
\end{definition}

In particular, two-term semorthogonal decompositions \(\mathscr{D}=\langle
\mathscr{D}_1,\mathscr{D}_2\rangle\) are determined by two additive subcategories
\(\mathscr{D}_1,\mathscr{D}_2\) that are closed unders shifts and such that
every \(E\in \mathscr{D}\) fits in an exact triangle \( E_2\to E\to
E_1\dashrightarrow\) for \(E_1\in \mathscr{D}_1\), \(E_2\in \mathscr{D}_2\). Finding
semi-orthogonal decompositions of \(\mathscr{D}_1\) and \(\mathscr{D}_2\) and
iterating gives semi-orthogonal decompositions of \(\mathscr{D}\) with
multiple terms.

Thus, for example, every partition of a full exceptional collection gives a
semi-orthogonal decomposition. If \(L\) is an exceptional object, then the
corresponding component \(\langle L \rangle\) is equivalent to the derived
category of a point (since every object has form \(L\otimes V^\bullet\) for a
complex of vector spaces \(V^\bullet\)).

\begin{example}
    The category \(\deri(\mathbb{P}^1)\) is ``built out of two points'', in the
    sense that there is a full exceptional collection
    \((\mathscr{O}_{\mathbb{P}^1},\mathscr{O}_{\mathbb{P}^1}(-1))\).  Indeed we
    have already seen the orthogonality condition, and if \(E\in
    \deri(\mathbb{P}^1)\) is indecomposable then it is either of the form
    \(\mathscr{O}_{\mathbb{P}^1}(k)[i]\) for some \(k,i\in \mathbb{Z}\) or of
    the form \(\mathscr{O}_x[i]\) for some skyscraper sheaf \(\mathscr{O}_x\)
    and \(i\in \mathbb{Z}\). The existence of the required exact triangle can
    then be deduced from triangles of the form
    \(\mathscr{O}_{\mathbb{P}^1}^{\oplus k-1}\to
    \mathscr{O}_{\mathbb{P}^1}(1)^{\oplus k} \to
    \mathscr{O}_{\mathbb{P}^1}(k)\dashrightarrow\) and \(
    \mathscr{O}_{\mathbb{P}^1}\to \mathscr{O}_x\to
    \mathscr{O}_{\mathbb{P}^1}(-1)[1]\dashrightarrow\).
\end{example}

\subsection{Admissible subcategories} Semi-orthogonal decompositions can be
determined from just one component, hence it suffices to look for certain
well-behaved subcategories.

\begin{definition}
    A full subcategory \(\mathscr{D}'\subset \mathscr{D}\) is \textit{left
    (right) admissible} if the inclusion functor has a left (resp.\ right)
    adjoint. We say \(\mathscr{D}'\) is \textit{admissible} if it is both left
    and right admissible.
\end{definition}

\begin{proposition}
    If \(\mathscr{D}=\langle \mathscr{D}_1,\mathscr{D}_2\rangle\) is a
    semi-orthogonal decomposition, then \(\mathscr{D}_1\) is left admissible and
    \(\mathscr{D}_2\) is right admissible. Conversely if \(\mathscr{D}_1\subset
    \mathscr{D}\) is a left admissible subcategory then \(\langle
    \mathscr{D}_1, \orth \mathscr{D}_1\rangle\) is a semi-orthogonal
    decomposition. Likewise for right admissible subcategories
    \(\mathscr{D}_2\subset \mathscr{D}\).

    \begin{proof}
        If \(\mathscr{D}=\langle \mathscr{D}_1,\mathscr{D}_2\rangle\) is a
        semi-orthogonal decomposition then every \(E\in \mathscr{D}\) sits in a
        unique(!) triangle \(E_2\to E\to E_1\dashrightarrow\) with \(E_i\in
        \mathscr{D}_i\). Then the required adjoint functors are given by
        \(E\mapsto E_i\). Conversely if the inclusion \(i:\mathscr{D}_2\to
        \mathscr{D}\) has a right adjoint \(i^!:\mathscr{D}\to \mathscr{D}_2\),
        then for any \(E\in \mathscr{D}\) we have the adjunction map \(E\to
        i\circ i^!E\) with cone \(E_1\). Remains to show \(E_1\in
        \mathscr{D}_2\orth\), which holds since for any \(F\in \mathscr{D}_2\)
        we have \[\Hom^\bullet(iF, E)=\Hom^\bullet(F,i^!E)=\Hom^\bullet(iF,
        i\circ i^!E)\] i.e.\ the term \(\Hom^\bullet(iF, E_1)\) must vanish in
        the long exact sequence associated to \(E\to i\circ i^!E\to
        E_1\dashrightarrow\).
    \end{proof}
\end{proposition}

\begin{remark}
    If \(\mathscr{D}\) has a Serre functor \(\mathbb{S}\) (e.g.\ if
    \(\mathscr{D}=\deri(X)\) for \(X\) smooth) then
    semi-orthogonal decompositions always come in pairs \(\mathscr{D}=\langle
    \mathscr{D}_1,\mathscr{D}_2 \rangle = \langle \mathbb{S}\mathscr{D}_2,
    \mathscr{D}_1 \rangle\). In particular a category is left admissible if and
    only if it is right admissible. Now if \(\mathscr{D}=\langle
    \mathscr{D}_1,\mathscr{D}_2,\mathscr{D}_3\rangle\) is a semi-orthogonal
    decomposition such that the subcategory \(\langle \mathscr{D}_2,
    \mathscr{D}_3\rangle\) has a Serre functor \(\mathbb{S}_{23}\), then we
    obtain a new semi-orthogonal decomposition \(\mathscr{D}=\langle
    \mathscr{D}_1, \mathbb{S}_{23}\mathscr{D}_3,\mathscr{D}_2\rangle\). This
    general process of \textit{mutation} can be seen as an action of a braid
    group on the semi-orthogonal decomposition.
\end{remark}

Where to find admissible subcategories? The following theorem gives one source.

\begin{theorem}
    If \(X,Y\) are smooth varieties and \(\Phi:\deri(X)\to \deri(Y)\) is fully faithful,
    then the image of \(\Phi\) is admissible.
\end{theorem}

Exceptional objects can be seen as a special instance of the above. Note that \(E\in \mathscr{D}\) is an exceptional object if and only if the map
\(\deri(\text{pt})\to \mathscr{D}\) given by \(V^\bullet \mapsto E\otimes
V^\bullet\) is fully faithful. The image is right admissible as \(F\mapsto
\Hom(E,F)\) is the right adjoint by the tensor--hom adjunction.

Thus exceptional object induces a semi-orthogonal decomposition
\(\mathscr{D}=\langle E\orth, E\rangle\). Iterating this construction, we can
induce semi-orthogonal decompositions from exceptional collections.

\begin{definition}
    If \(E_1,...,E_m\) is an exceptional collection in \(\deri(X)\) then the
    \textit{Kuznetsov component} of \(\deri(X)\) with respect to this exceptional
    collection is \(\text{Ku}(X)=E_1\orth \cap E_2\orth \cap ... \cap
    E_m\orth\). Thus \(\deri(X)=\langle \text{Ku}(X),E_1,...,E_m\rangle\) is a
    semi-orthogonal decomposition.
\end{definition}

While the definition of a Kuznetsov component depends on the choice of an
exceptional collection, such a choice is often canonically made. We will see
some examples.


\subsection{Semi-orthogonal decompositions on families} Take a quadric
\(Q\subset \mathbb{P}^3\). If \(Q\) is smooth then it is isomorphic to
\(\mathbb{P}^1\times \mathbb{P}^1\), so there is a full exceptional collection
\((\mathscr{O}_Q(-1,0), \mathscr{O}_Q(0,-1), \mathscr{O}_Q, \mathscr{O}_Q(1,1))\). Of these, the latter two are canonical (they come from the structure sheaf and the canonical bundle) so we choose those as
our exceptional objects and hence we have the Kuznetsov component
\[\text{Ku}(Q)=\langle \mathscr{O}_Q(-1,0),\mathscr{O}_Q(0,-1)
\rangle=\deri(\text{pt})^{\oplus 2}.\]

If \(Q\) is nodal singular (i.e.\ given by \(\mathbb{V}(xy-z^2)\subset
\mathbb{P}^3\)), then the Kuznetsov component is generated by the ideal sheaf of
the line \(\ell=\mathbb{V}(x,z)\) so we have
\(\text{Ku}(Q)=\deri(k[\epsilon]/(\epsilon^2))\).  The intuition is that as the
quadric surface degenerates from smooth to singular, the Kuznetsov component
undergoes a similar transition.  To formalise this, we will study families of
varieties.

\begin{definition}
    Let \(\pi:X\to B\) be a Gorenstein morphism. A \textit{relative
    exceptional object} is a perfect complex \(E\in \deri(X)\) such that
    \(\sheafHom_B^\bullet(E,E)\coloneq \pi_\ast \sheafHom^\bullet(E,E)\cong
    \mathscr{O}_B[0]\). A \textit{relative exceptional collection} is a sequence
    of relative exceptional objects \(E_1,...,E_m\) with
    \(\sheafHom_B(E_i,E_j)=0\) for \(i>j\).
\end{definition}

The Gorenstein assumption guarantees the Serre functor takes bounded complexes
to bounded complexes, and the assertion that \(E\) is perfect (i.e.\
quasi-isomorphic to a bounded complex of vector bundles) guarantees that it is
well-behaved under tensor products. It follows that the map \(\pi^\ast(-)\otimes
E:\deri(B)\to \deri(X)\) is fully faithful with admissible image, so
we can induce semi-orthogonal decompositions as before.

Going back to quadrics, let \(\pi:\mathcal{Q}\to B\) be a family of quadric
surfaces with \(\dim B = 1\), and \(\mathcal{Q}, B\) smooth. It follows that the
fibers are at worst nodal. Then \((\mathscr{O}_{\mathcal{Q}},
\mathscr{O}_{\mathcal{Q}}(1))\) is a relative exceptional collection giving us a
semi-orthogonal decomposition
\[
    \deri(\mathcal{Q})=\left\langle \text{Ku}(\mathcal{Q}), \pi^\ast \deri(B),
    \pi^\ast \deri(B)\otimes \mathscr{O}_{\mathcal{Q}}(1) \right\rangle.
\]
To understand the Kuznetsov component better, consider the 2:1 cover associated
to \(\pi\) obtained as follows. We have the relative Hilbert scheme
\(\text{Hilb}_\text{lines}(\mathcal{Q})\to B\), with Stein factorisation
\(\text{Hilb}_\text{lines}(\mathcal{Q})\to \tilde B\to B\). Thus the map
\(\text{Hilb}_\text{lines}(\mathcal{Q})\to \tilde B\) is proper with connected
fibers, while the map \(\tilde B\to B\) is a \(2:1\) cover branched at the
discriminant locus in \(B\).

We can think of \(\tilde B\) is the relative moduli space of torsion-free
sheaves \(F\) in the fibers of \(\pi\) with an equality of Hilbert polynomials \( p_H(F)=p_H(\mathscr{O}_{\mathbb{P}^1\times \mathbb{P}^1}(-1,0)).\)
As a result there is a universal family in \(\Coh(\tilde B \times
\mathcal{Q})\), giving a Fourier--Mukai transform \(\Phi:\deri(\tilde B)\to
\deri(\mathcal{Q})\). It can be shown that this is fully faithful with image
\(\text{Ku}(\mathcal{Q})\), so that the semi-orthogonal decomposition is
\[
    \deri(\mathcal{Q})=\langle \deri(\tilde B), \pi^\ast \deri(B), \pi^\ast
    \deri(B)\otimes \mathscr{O}_{\mathcal{Q}}(1) \rangle.
\]
Since relative exceptional collections are well-behaved under pullbacks, this
explains the relation between the Kuznetsov components of smooth and nodal
quadrics-- if the fiber of \(\mathcal{Q}\) over \(b\in B\) is smooth, the
moduli space \(\tilde B\) has two points corresponding to \(\mathscr{O}(-1,0)\)
and \(\mathscr{O}(0,-1)\). On the other hand \(\tilde B\to B\) is branched
precisely over the \(b\in B\) corresponding to nodal quadrics, where the fiber
of \(\tilde B\) looks like \(\Spec k[\epsilon]/(\epsilon^2)\).

\subsection{Decompositions via birational transformations} Orlov's blow-up
formula explains how Kuznetsov components transform under birational maps-- in
general blowing up along a codimension \(r\) locus \(Z\into X\) adds \(r-1\)
copies of \(\deri(Z)\) twisted by some exceptional bundles to \(\deri(\text{Bl}_ZX)\).

More specifically if \(X\) is a 3-fold with a smooth point \(x\in X\), then the
blow-up \(p:\text{Bl}_xX\to X\) has exceptional fiber \(E\cong \mathbb{P}^2\)
and we have \(\deri(\text{Bl}_xX)= \langle \mathscr{O}_E(-2), \mathscr{O}_E(-1),
p^\ast \deri(X)\rangle\). In particular, the Kuznetsov component of \(X\)  is
``essentially undamaged'', it simply picks up two points. On the other hand if
\(C\subset X\) is a smooth curve then we have \(\deri(\text{Bl}_CX)=\langle
\deri(C), p^\ast \deri(X) \rangle\) i.e.\ the Kuznetsov component  picks up
\(\deri(C)\). Note the striking similarity with intermediate Jacobians!

To see a concrete example, fix two quadrics \(Q_1,Q_2\subset \mathbb{P}^3\) and
consider the associated pencil. This gives a family \(\mathcal{Q}\) over
\(B=\mathbb{P}^1\), which can be realised as a blow-up \(p:\mathcal{Q} \to
\mathbb{P}^3\) along the elliptic curve \(C=Q_1\cap Q_2\) giving us
\[
    \deri(\mathcal{Q})=\langle \deri(C), p^\ast \mathscr{O}_{\mathbb{P}^3},
    p^\ast \mathscr{O}_{\mathbb{P}^3}(1), p^\ast
    \mathscr{O}_{\mathbb{P}^3}(2),p^\ast \mathscr{O}_{\mathbb{P}^3}(3)  \rangle.
\]

How does this relate to the semi-orthogonal decomposition computed previously?
Note the map \(\pi:\mathcal{Q}\to B\) has singular fibers at four points in
\(B\), so the \(2:1\) branched cover \(\tilde B\to B\) is an elliptic curve.
In particular, both the semi-orthogonal decompositions are made of the derived
category of an elliptic curve and four exceptional objects. These are in fact
related and \(\tilde B = C\), but the proof is somewhat ad hoc.

\begin{remark}
    We just saw that it is possible to have multiple semi-orthogonal
    decompositions that are not related in an obvious way (eg.\ by mutation).
    This is quite typical! There is however a conjectural proposal by
    Halpern--Leistner relating semi-orthogonal decompositions that arise
    `geometrically'.
\end{remark}

\begin{example}
    Playing the above game in \(\mathbb{P}^5\), consider the pencil associated
    to two general quadrics \(Q_1,Q_2\). The resulting \((2,2)\)-complete
    intersection \(Q_1\cap Q_2\) is the index \(2\) Fano 3-fold \(Y_4\), and
    the pencil can be realised as a blow-up
    \(\mathcal{Q}=\text{Bl}_{Y_4}(\mathbb{P}^5)\to \mathbb{P}^1\). Thus we again
    have two ways to compute \(\deri(\mathcal{Q})\)-- since the blow-up is along
    a codimension \(2\) locus, Orlov's blow-up formula tells us
    \[\deri(\mathcal{Q})=\langle \deri(Y_4), p^\ast \deri(\mathbb{P}^5)\rangle.\]
    On the other hand, the pencil \(\pi:\mathcal{Q}\to \mathbb{P}^1\)
    has singular fibers over six points, so that the
    associated double cover \(\tilde B\) is a genus \(2\)
    curve. A general smooth quadric \(Q\) in the pencil has an exceptional
    collection \(\mathscr{O}_Q, \mathscr{O}_Q(1),\mathscr{O}_Q(2),
    \mathscr{O}_Q(3)\) and a Kuznetsov component coming from two spinor bundles,
    which we can use to find a relative exceptional collection on \(\tilde{Q}\).
    This, together with the derived category of \(\tilde{B}\) give us a
    semi-orthogonal decomposition
    \[
        \deri(\tilde Q) = \langle \deri(\tilde B) , \pi^\ast \deri(\mathbb{P}^1),
        ..., \pi^\ast \deri(\mathbb{P}^1)\otimes \mathscr{O}_{\tilde
        Q}(3)\rangle.
    \]
    Note that one semi-orthogonal decomposition has six exceptional objects
    (coming from \(\deri(\mathbb{P}^5)\)), while the other has eight! This was
    explained by Bondal and Orlov who recovered the two exceptional objects by
    computing \(\deri(Y_4)=\langle \deri(\tilde B), \mathscr{O}_{Y_4},
    \mathscr{O}_{Y_4}(1) \rangle.\) In fact such a formula holds for all prime
    Fano 3-folds of index \(2\).
\end{example}

\subsection{Mukai bundles} We discuss a powerful result on how to find
exceptional objects.

\begin{theorem}[(Mukai)]
    If \(X\) is a prime Fano 3-fold of index \(1\) and genus \(g\geq 6\), then
    for \(r=2,3\) and \(s\geq r\) satisfying \(rs=g\) there is a unique vector
    bundle \(\mathcal{E}_r\) of rank \(r\) and \(c_1(\mathcal{E}_r)=-H\) that is
    exceptional. Moreover, it is stable and satisfies
    \(\HH^\bullet (X,\mathcal{E}_r)=0\). Its dual \(\mathcal{E}_r\dual\) is
    globally generated by \(r+s\) sections, giving a map \(X\to
    \textnormal{Gr}(r,r+s)\). This map is an embedding unless \(g=6\) and \(X\)
    is ``special'', i.e.\ a double cover of \(Y_5\).
\end{theorem}

It is hard to overstate the importance of this theorem. For instance, it leads
to the classification of prime Fano 3-folds of index \(1\). In particular they
are complete intersections for \(g\leq 5\). For \(g=6\), they are either
special or \((2,1,1)\)-complete intersections in \(\text{Gr}(2,5)\). For
\(g=8\) they are linear sections of \(\text{Gr}(2,6)\) and for genus
\(10\) they are obtained as the intersection of a linear section of
\(\text{Gr}(2,6)\) with the zero section of a vector bundle.

As a second application, we have a canonical choice of exceptional objects
\(\mathcal{E}_r, \mathscr{O}_X\) so the Kuznetsov component of \(X\) is defined
with respect to these, and we have \(\deri(X)=\langle\text{Ku}(X), \mathcal{E}_r,
\mathscr{O}_X \rangle\).

\textbf{But there's a problem.} There is no complete and correct proof in the
literature, so a proof of this result is a work in progress
[Bayer--Kuznetsov--Macri].

\subsection{Perry's categorical intermediate Jacobian} Kuznetsov components and
intermediate Jacobians share many similar properties-- an attempt to explain
this is a categorification of Jacobians. An ingredient in this
categorical Jacobian are Blanc's topological \(K\)-theories \(K_0^\text{top}\mathscr{D},
K_1^\text{top}\mathscr{D}\) defined for a suitably enriched
\(\mathbb{C}\)-linear category \(\mathscr{D}\). When
\(\mathscr{D}=\deri(X)\) this coincides with the topological \(K\)-theory
\(K_i^\text{top}\mathscr{D}=K_i^\text{top}X\) and we have
\(K^\text{top}_{0,1}\mathscr{D}\otimes
\mathbb{Q}=\HH^{\text{even},\text{odd}}(X,\mathbb{Q})\).

The construction is additive on semi-orthogonal decompositions, and has an Euler
pairing. There are two Chern characters-- one from the usual Grothendieck group
\(K_0\mathscr{D}\to K_0^\text{top}\mathscr{D}\), and another to the Hochschild
cohomology which has a canonical filtration. This can be used to show the
following result.

\begin{theorem}[(Perry)]
    Assume \(\mathscr{D}\subset \deri(X)\) is admissible. Then
    \(K^\textnormal{top}_{0,1}(\mathscr{D}\) carries a natural integral polarised
    Hodge structure compatible with the above properties. Moreover if
    \(n=\dim X\) is odd and
    \(\HH^\textnormal{odd}(X,\mathbb{Z})=\HH^n(X,\mathbb{Z})\), then
    \(\HH^n(X,\mathbb{C})\cong K_1^\textnormal{top}\deri(X)\) as integral Hodge
    structures. If additionally \(\HH^{p,n-p}(X)=0\) except for
    \(p=\frac{n\pm 1}{2}\), then \(\HH^n(X)\) coincides with
    \(\IJac X\).
\end{theorem}

\begin{corollary}
    In the above situation, a semiorthogonal decomposition \(\deri(X)=\langle
    \mathscr{D}_1,...,\mathscr{D}_n \rangle\) gives a decomposition
    \(\IJac X=\IJac \mathscr{D}_1\times ... \times
    \IJac \mathscr{D}_n\) where the Jacobians \(\IJac \mathscr{D}_i\)
    are computed from the weight \(1\) intermediate Hodge structure on
    \(K_1^\text{top}\mathscr{D}_i\).
\end{corollary}

\section{Stability conditions and moduli spaces}
We will now see how to elucidate Kuznetsov components of Fano 3-folds by
studying moduli spaces of Bridgeland-semistable objects.

\begin{definition}
    A \textit{stability condition} on an Abelian category \(\mathscr{A}\) is an
    additive map \(Z:K_0\mathscr{A} \to \mathbb{Z}\) (called the \textit{central
    charge}) that satisfies \(Z(E)\in \mathbb{H}_-\) for any \(0\neq E\in
    \mathscr{A}\) and has the Harder--Narasimhan property.
\end{definition}

We explain what the Harder--Narasimhan property is. Note since \(Z\) sends the
class of a non-zero object into the semi-closed upper half plane, the
\textit{phase} of an object is a well defined real number in \((0,1]\). We say
\(0\neq E\in \mathscr{A}\) is \(Z\)-semistable if all non-zero sub-objects of
\(E\) have smaller phase, or equivalently a smaller slope (where the slope of
\(E\) is \(\mu(E)=-\frac{\text{Re}ZE}{\text{Im}ZE}\)). Then \(Z\) satisfies the
Harder--Narasimhan property if every \(0\neq E\in \mathscr{E}\) admits a
filtration \(0=E_0\subset E_1\subset ... \subset E_n = E\) such that all factors
\(E_i/E_{i-1}\) are semistable and satisfy
\(\mu(E_1/E_0)>\mu(E_2/E_1)>...>\mu(E_n/E_{n-1})\).

\begin{example}
    If \(\mathscr{A}=\Coh C\) for a smooth projective curve, then
    \(Z=i\cdot \text{rank}-\text{deg}\) gives a stability condition.
\end{example}

We also define a \textit{weak stability condition}, in which we allow non-zero
objects \(E\) to have \(Z(E)=0\) (in which case the slope is \(+\infty\)).

\begin{example}
    If \(S\) is a smooth projective surface and \(H=[C]\) is a polarisation,
    then \(Z=i\cdot \text{rank}-H\cdot c_1\) (i.e.\ \(Z(E)= i\cdot
    \text{rank}(E|_C)-\text{deg}(E|_C)\)) is a weak stability function.
\end{example}

To extend this to triangulated categories, we define first the `heart of a
bounded t-structure'. The prime example is an Abelian category \(\mathscr{A}\)
as a subcategory of \(\deri^b \mathscr{A}\).

\begin{definition}
    If \(\mathscr{D}\) is a triangulated category, then a \textit{heart} is a
    full additive subcategory \(\mathscr{A}\subset \mathscr{D}\) such that
    \(\mathscr{A}[k_1]\subset \orth\mathscr{A}[k_2]\)  whenever \(k_1>k_2\),
    and every \(0\neq E\in \mathscr{D}\) admits a filtration (i.e.\ a sequence
    of maps)
    \[
        0=E_0\to E_1\to ... \to E_n = E
    \]
    with \(\cone(E_{i-1}\to E_i)\in \mathscr{A}[k_i]\) for \(k_1>k_2>...\).
\end{definition}

To draw the analogy with stability functions on Abelian categories, observe
that if \(E,F\in \mathscr{A}\) are semistable objects with slopes
\(\mu(E)>\mu(F)\), then \(\Hom(E,F)=0\). Thus \(\mathscr{A}\) is filtrated by
`slices' of semistable objects much like how \(\mathscr{D}\) is filtrated by
shifts of \(\mathscr{A}\). A stability condition combines these two ideas.

\begin{definition}
    A \textit{stability condition} on \(\mathscr{D}\) is given by the heart of a
    bounded t-structure \(\mathscr{A}\) and a stability function on
    \(\mathscr{A}\). A weak stability condition is defined analogously.
\end{definition}

\subsection{Tilting of hearts} How to construct hearts? We have already seen
\(\mathscr{A}\) is a heart of \(\deri^b\mathscr{A}\). If \(\Phi\) is an
autoequivalence of \(\mathscr{D}\), then applying \(\Phi\) to existing hearts
also gives new hearts. But we need a more systematic way to get new hearts, and
this can be done by perturbing existing ones.

\begin{definition}
    A torsion pair \((T,F)\) in an Abelian category \(\mathscr{A}\)  is a pair
    of full additive subcategories such that \(T\subset \orth F\) and every
    \(E\in \mathscr{A}\) sits in a short exact sequence \(0\to E_T\to E\to
    E_F\to 0\) with \(E_T\in T\), \(E_F\in F\).
\end{definition}

The prime example is that of torsion sheaves and torsion-free sheaves on a
smooth variety.

Torsion pairs can be thought of as a coarse filtration of \(\mathscr{A}\) into
two blocks. Thus for instance if \(\mathscr{A}\) has a stability function, then
for any \(\mu\in (0,1] \) we have an induced torsion pair
\(T=\mathscr{A}^{>\mu}, F=\mathscr{A}^{\leq \mu}\) where the torsion part
contains all objects whose semistable factors have slope \(>\mu\).

\begin{proposition}
    If \(\mathscr{A}\) is the heart of a bounded t-structure and \((T,F)\) is a
    torsion pair on \(\mathscr{A}\), then the extension closure
    \(\mathscr{A}^\sharp=\langle T, F[1] \rangle\) is the heart of a bounded
    t-structure. Hearts that arise in this way are called \emph{tilts} of
    \(\mathscr{A}\).
\end{proposition}

Concretely, if \(\mathscr{A}\) is the natural heart of \(\deri^b\mathscr{A}\)
then \(\mathscr{A}^\sharp\) contains precisely those complexes \(E\) which have
\(\HH^{-1}(E)\in F\), \(\HH^0(E)\in T\), and \(\HH^i(E)=0\) otherwise. Such
\(E\) can always be represented by a 2-term complex \(E^{-1}\xrightarrow{d}E^0\)
with \(\ker d\in F\), \(\coker d\in T\).

\subsection{Constructing hearts in Kuznetsov components} In general, the process
of taking semi-orthogonal decompositions is not compatible with taking hearts of
bounded t-structures. However, if we are dealing with a single exceptional
object then hearts descend to Kuznetsov components under reasonable assumptions.

\begin{proposition}
    Suppose \(\mathscr{D}\) has an exceptional object \(E\) that lies in a heart
    \(\mathscr{A}\). Moreover, suppose \(\Hom^q(E,F)=0\) for all \(F\in
    \mathscr{A}\) unless \(q=0,1\). Then \(\mathscr{A}\cap E\orth\) is the
    heart of a bounded structure in \(E\orth\).

    \begin{proof}
        For \(F\in E\orth\), we wish to show that each cohomology object with
        respect to \(\mathscr{A}\) is also in \(E\orth\). Now there is a
        spectral sequence \(E_2^{p,q}=\Hom^q(E, \HH_{\mathscr{A}}^pF)\Rightarrow
        \Hom^{p+q}(E,F)=0\) which shows that \(\Hom^q(E,
        \HH_{\mathscr{A}}^pF)=0\) for all \(p,q\) as required.
    \end{proof}
\end{proposition}

 Note if \(\mathscr{D}\) has an exceptional object \(E\) that lies in a heart
 \(\mathscr{A}\) such that \(\mathbb{S}E\in \mathscr{A}[1]\) under the Serre
 functor, then the hypotheses of the above theorem are satisfied. Indeed,
 for \(F\in \mathscr{A}\) and \(q\geq 2\) we have \(\Hom(E,F[q])=\Hom(F[q],
 \mathbb{S}E)=0\).

\begin{proposition}
    In the above setting, if \(Z\) gives a stability condition on
    \(\mathscr{A}\) with \(\textnormal{Im}Z(E)>0\) then \(Z|_{\mathscr{A}\cap
    E\orth}\) is a stability condition on \(\mathscr{A}\cap E\orth\).
\end{proposition}

\begin{example}
    For a del Pezzo surface \(S\) with polarisation \(H\), one can consider
    the corresponding weak stability condition and note that there is an
    exceptional collection \(E_1,...,E_m\) such that for all \(i,j\) we have
    \(\mu_H(E_i)>\mu_H(E_j)-H\cdot (-K_S)\). If all \(E_i\) are semistable (this
    holds for the canonical polarisation, for example) then the Kuznetsov
    component \(E_1\orth \cap ...
    \cap E_m\orth\) has a bounded t-structure obtained as follows. Take \(\mu\in
    \mathbb{R}\) such that \(\mu_H(E_i)-H(-K_S)\leq \mu<\mu_H(E_i)\) for all
    \(i\), and consider the torsion pair \((\Coh^{{>\mu}}S, \Coh^{\leq \mu}S)\)
    on \(\Coh S\).  The tilt \(\Coh^\mu S\) contains all \(E_i\), and we have
    \(\mathbb{S}E_i = E_i(K_S)[2]\in \Coh^\mu [1]\) by a simple slope
    computation. Thus \(\Coh^\mu S \cap E_1\orth \cap ... \cap E_m\orth\) is the
    heart of a t-structure.

    As a corollary, we see that such collections on del Pezzo surfaces have no
    phantoms. Indeed if \(K_0(E_1\orth \cap ... \cap E_m\orth)\) is torsion, then
    there are no stability conditions and hence we must have \(E_1\orth \cap ...
    \cap E_m\orth = 0\) (otherwise we could restrict stability conditions from
    \(\Coh S\).)
\end{example}

The above ideas can be used to show the following.

\begin{theorem}
    For any Fano 3-fold of index \(\leq 2\), there is a stability condition on
    the Kuznetsov component.
\end{theorem}

It was shown by Pertrusi and Yang that these stability conditions are often
invariant under the Serre functor on the Kuznetsov component. Moreover often
there is a unique stability condition with such invariance property. This leads
to Torelli--type theorems.

\begin{example}[The cubic 3-fold] Recall that \(\deri(Y_3)=\langle \text{Ku}(Y_3),
    \mathscr{O}_{Y_3}, \mathscr{O}_{Y_3}(1)\rangle\) so let \(i\) be the
    inclusion of the Kuznetsov component. Consider the pulled back
    skyscraper \(i^\ast\mathscr{O}_y\) for \(y\in Y_3\), which sits in an exact
    sequence \(I_y(1)\into \mathscr{O}_{Y_3}(1)\onto \mathscr{O}_y \). But there
    is a surjection \(\mathscr{O}^{\oplus 4}_{Y_3}\onto
    I_y(1)\) with kernel \(K_y\), and we see that \(i^\ast
    \mathscr{O}_y=K_y(2)\).

    It is then a theorem that the sheaves \(K_y\) belong to a four dimensional
    moduli space of stable sheaves \(\Moduli(v)\), where \(v\) is the Chern
    character. In particular there is an embedding \(Y_3\into \Moduli(v)\)
    given by \(y\mapsto K_y\). Sheaves \(E\) in the complement
    \(\Moduli(v)\setminus Y_3\) arise from the following construction: pick
    a cubic surface \(S_3\subset Y_3\). This has seventy two maps to
    \(\mathbb{P}^2\), and for any \(\pi:S_3\to \mathbb{P}^2\) we have that
    \(\pi^\ast\mathscr{O}(1)\) is a divisor of the form \(H+L_1-L_0\). Here
    \(L_1-L_0\) is a root of the associated \(E_6\) root system. More
    importantly, there is a short exact sequence \(E\into
    \mathscr{O}_{Y_3}^{\oplus 3}\onto \mathscr{O}_{S_3}(H+L_1-L_0)\) which
    shows \(E\) is locally free. It follows that \(Y_3\) is precisely the locus
    of non locally--free sheaves in \(\Moduli(v)\).

    But the moduli space also has an Abel--Jacobi map \(\Moduli(v)\to
    \IJac Y_3=\text{CH}_1(Y_3)\) and under this map, we have
    \(K_y\mapsto 0\) while other objects \(E\) get mapped to \([L_1]-[L_0]\).
    In their analysis of the theta divisor \(\Theta\subset \IJac Y_3\),
    Clemens and Griffiths show that it is given by classes of the form
    \([L_1]-[L_0]\). Thus \(\Moduli (v)\) maps surjectively onto \(\Theta\)
    and \(Y_3\subset \Moduli(v)\) lies in the kernel. Computing normal bundle
    sheaves, they show that \(0\in \Theta\) coincides with \(0\in
    \text{Cone}(Y_3)\) formally locally. The map is an isomorphism away from
    \(Y_3\in \Moduli(v)\), \(0\in \Theta\). Thus \(\Theta\) is smooth away
    from \(0\) and \(Y_3\) is the exceptional locus of the blow--up, giving
    another way to derive the Torelli theorem for cubic 3-folds.

    What does this have to do with Kuznetsov components? In this case we can
    show \(\Moduli(v)\) is also the moduli space of all stable objects in
    \(\text{Ku}(Y_3)\). As a corollary, we get a very short proof of the theorem
    that \(\text{Ku}(Y_3)\) determines \(Y_3\). Indeed there is a unique
    stability condition invariant under the Serre functor, and this gives a four
    dimensional moduli space \(\Moduli(v)\). The 3-fold \(Y_3\) lies in this
    space as the union of all rational curves.
\end{example}
